\documentclass[../../../include/open-logic-section]{subfiles}

\begin{document}

\olfileid[fr]{sth}{choice}{hartogs}
\olsection{Comparabilité et lemme de Hartogs}

Voilà pour les avantages. Passons aux difficultés :
sans l'axiome du choix, les choses se
\emph{compliquent}. Pour comprendre pourquoi,
considérons un résultat élégant dû à
\cite{Hartogs1915} :

\begin{lem}[\emph{dans $\ZF$}]\ollabel{HartogsLemma}
Pour tout ensemble $A$, il existe un ordinal
$\alpha$ tel que $\cardnless{\alpha}{A}$.
\end{lem}

\begin{proof}
Si $B \subseteq A$ et $R \subseteq B^2$,
alors $\tuple{B, R} \subseteq
V_{\setrank{A}+4}$ d'après le
\olref[ord-arithmetic][using-addition]{rankcomputation}.
Par séparation, considérons donc :
\[
	C = \Setabs{\tuple{B, R} \in V_{\setrank{A}+5}}
	{B\subseteq A
	\text{ et $\tuple{B, R}$ est un bon ordre}}
\]
Par remplacement et par le
\olref[ordinals][ordtype]{thmOrdinalRepresentation},
formons l'ensemble :
\[
	\alpha =
	\Setabs{\ordtype{B, R}}{\tuple{B, R} \in C}.
\]
D'après
le \olref[ordinals][basic]{corordtransitiveord},
$\alpha$ est un ordinal, car c'est un
ensemble transitif d'ordinaux. En effet,
si $\gamma \in \beta \in \alpha$, alors
$\beta = \ordtype{B, R}$ pour un certain
$B \subseteq A$ ; de plus,
$\gamma = \ordtype{B_b, R_b}$ pour un
$b \in B$ par le
\olref[ordinals][iso]{wellordinitialsegment},
si bien que $\gamma \in \alpha$.\footnote{
Note de l'édition française : l'original
écrit ici $B\subseteq R$, alors que
la définition de $C$ impose $B\subseteq A$.}

Pour raisonner par l'absurde, supposons
qu'il existe !!a{injection}
$f \colon \alpha \to A$. Posons :
\begin{align*}
	B &= \ran{f}\\
	R &= \Setabs{\tuple{f(\xi), f(\eta)}}
	{\xi,\eta\in\alpha \text{ et } \xi \in \eta}.
\end{align*}
\footnote{Note de l'édition française :
l'original réutilise $\alpha$ dans la
définition de $R$, alors que cette lettre
désigne aussi le domaine de $f$.
Nous explicitons les indices liés
$\xi,\eta\in\alpha$ pour lever l'ambiguïté
et transportons leur ordre par $f$.}
Il est clair que
$\alpha = \ordtype{B, R}$ et que
$\tuple{B, R} \in C$. Donc
$\alpha \in \alpha$, contradiction.
\end{proof}

Il en découle un résultat profond :

\begin{thm}[\emph{dans $\ZF$}]
Les assertions suivantes sont équivalentes :
\begin{enumerate}
	\item\ollabel{equivwo}
	L'axiome du bon ordre ;
	\item\ollabel{equivcompare}
	Pour tous ensembles $A$ et $B$,
	$\cardle{A}{B}$ ou $\cardle{B}{A}$.
\end{enumerate}
\end{thm}

\begin{proof}
\emph{\olref{equivwo} $\Rightarrow$
\olref{equivcompare}.}
Fixons $A$ et $B$. Par \olref{equivwo},
il existe des bons ordres
$\tuple{A, R}$ et $\tuple{B, S}$.
D'après le
\olref[ordinals][ordtype]{thmOrdinalRepresentation},
soient $f \colon \alpha \to \tuple{A, R}$
et $g \colon \beta \to \tuple{B, S}$
des isomorphismes. D'après la
\olref[sth][ordinals][basic]{ordinalsaresubsets},
soit $\alpha \subseteq \beta$, soit
$\beta \subseteq \alpha$.
Dans le premier cas,
$\comp{f^{-1}}{g} \colon A \to B$
est !!a{injection}, et donc
$\cardle{A}{B}$ ; dans le second,
on obtient de même $\cardle{B}{A}$.

\emph{\olref{equivcompare} $\Rightarrow$
\olref{equivwo}.}
Fixons $A$. D'après le \olref{HartogsLemma},
il existe un ordinal $\beta$ tel que
$\cardnless{\beta}{A}$.
Par \olref{equivcompare}, nous avons
donc $\cardle{A}{\beta}$.
Il existe ainsi !!{injection}
$f \colon A \to \beta$, qui permet
de bien ordonner les éléments de $A$
au moyen de la relation
$\Setabs{\tuple{a, b} \in A \times A}
{f(a) \in f(b)}$.
\end{proof}
\noindent
Conséquence immédiate : si l'axiome
du bon ordre est faux, certains ensembles
sont \emph{littéralement incomparables}
quant à leur taille. Dans ce cas,
l'arithmétique cardinale transfinie
se complique. Par exemple, nous devons
abandonner l'idée que, si $A$ et $B$
sont infinis, alors
$\cardeq{A \disjointsum B}{M}$
et $\cardeq{A \times B}{M}$,
où $M$ est le plus grand des deux
ensembles $A$ et $B$
(voir le \olref[card-arithmetic][simp]{cardplustimesmax}).\footnote{
Note de l'édition française : la formule
de l'original contient une expression
$\cardeq{\cardeq{A \disjointsum B}{A \times B}}{M}$
mal formée : elle traite une assertion
d'équipotence comme un ensemble.
Nous formulons séparément les deux égalités
de taille visées par le résultat cité.}
Le problème est simple : si nous ne pouvons
pas \emph{comparer} la taille de $A$ et de $B$,
demander lequel est le plus grand n'a pas
de sens.

\end{document}
